
%%%%%%%%%%%%%%%%%%%%%%%%%%%%%%%%%%%%%%%%%%%%%%%%%%%%%%%%%%%%%%%%%%%%%%%%%%%%%%%%

\part{Quantum circuits}\label{part:circuits}

\begin{frame}
  \partpage
\end{frame}

%%%%%%%%%%%%%%%%%%%%%%%%%%%%%%%%%%%%%%%%%%%%%%%%%%%%%%%%%%%%%%%%%%%%%%%%%%%%%%%%

\begin{frame}{Quantum circuit model}

Quantum circuits are generalizations of Boolean circuits

\bigskip
input \hspace{1.5cm} transformation \hspace{1.0cm} output (probabilistic)

\bigskip
\hspace{0.75cm}
\Qcircuit @C=0.7em @R=0.7em {
\lstick{\ket{1}} & \gate{U_{1\,}}     & \multigate{1}{U_3} & \qw                & \multigate{1}{U_6} & \meter & \rstick{0} \\
\lstick{\ket{1}} & \multigate{1}{U_2} & \ghost{U_3}        & \gate{U_5}         & \ghost{U_6}        & \meter & \rstick{1} \\
\lstick{\ket{0}} & \ghost{U_2}        & \qw                & \multigate{1}{U_6} & \qw                & \meter & \rstick{0} \\
\lstick{\ket{1}} & \multigate{1}{U_3} & \gate{U_4}         & \ghost{U_5}        & \multigate{1}{U_7} & \meter & \rstick{1} \\
\lstick{\ket{0}} & \ghost{U_3}        & \qw                & \qw                & \ghost{U_7}        & \meter & \rstick{1}
}

\end{frame}

%%%%%%%%%%%%%%%%%%%%%%%%%%%%%%%%%%%%%%%%%%%%%%%%%%%%%%%%%%%%%%%

%% Classical Bit

%%%%%%%%%%%%%%%%%%%%%%%%%%%%%%%%%%%%%%%%%%%%%%%%%%%%%%%%%%%%%%%

\begin{frame}{Classical bit}

Classical bit (bit):\, $\mathbb{B}:=\{0,1\}$

\bigskip
\begin{itemize}
\item Basis state: either $0$ or $1$

\bigskip
\item General state: a probability distribution $p=(p_0,p_1)$ on $\mathbb{B}$
\end{itemize}

\end{frame}

%%%%%%%%%%%%%%%%%%%%%%%%%%%%%%%%%%%%%%%%%%%%%%%%%%%%%%%%%%%%%%%

%% Classical Register

%%%%%%%%%%%%%%%%%%%%%%%%%%%%%%%%%%%%%%%%%%%%%%%%%%%%%%%%%%%%%%%

\begin{frame}{Classical register}

Classical register: $\mathbb{B}^n := \underbrace{\mathbb{B} \times \mathbb{B} \times \ldots \times \mathbb{B}}_{n}$

\bigskip
\begin{itemize}
\item Basis state: a binary string $x\in\mathbb{B}^n$

\bigskip
\item General state: a probability distribution $p=(p_x \,:\, x\in\mathbb{B}^n)$ on $\mathbb{B}^n$ (written as a column vector)

\bigskip

\bigskip
Remark: Note that $p$ is a vector with positive entries that is normalized with respect to the $\ell_1$-norm (the sum of the absolute values of the entries)
\end{itemize}

\end{frame}

%%%%%%%%%%%%%%%%%%%%%%%%%%%%%%%%%%%%%%%%%%%%%%%%%%%%%%%%%%%%%%%

%% Classical Transformations

%%%%%%%%%%%%%%%%%%%%%%%%%%%%%%%%%%%%%%%%%%%%%%%%%%%%%%%%%%%%%%%

\begin{frame}{Classical transformation}

Transformations on the classical register $\mathbb{B}$ are described by stochastic matrices

\bigskip
Stochastic matrices preserve the $\ell_1$-norm, i.e., probability distributions are mapped on probability distributions

\bigskip
Let $p$ be the state of the register.  The state after the transformation $P$ is given by the matrix-vector-product 

\[
P p
\]

\end{frame}

%%%%%%%%%%%%%%%%%%%%%%%%%%%%%%%%%%%%%%%%%%%%%%%%%%%%%%%%%%%%%%%

%% QUBIT

%%%%%%%%%%%%%%%%%%%%%%%%%%%%%%%%%%%%%%%%%%%%%%%%%%%%%%%%%%%%%%%

\begin{frame}{Qubit}

Quantum bit (qubit):\, two-dimensional complex Hilbert space $\mathbb{C}^2$\,

\bigskip
\begin{itemize}
\item Computational basis states (classical states):

\bigskip
\[
|0\>:=\left(\begin{array}{c}1\\0\end{array}\right)
\quad\quad\mbox{and}\quad\quad
|1\>:=\left(\begin{array}{c}0\\1\end{array}\right)
\]

\bigskip
\item General states: \textcolor{blue}{superpositions}

\bigskip
\[
|\psi\> = 
\left(
\begin{array}{c}
\alpha_0 \\
\alpha_1
\end{array}
\right)=
\alpha_0 |0\> + \alpha_1 |1\>\,,\quad |\alpha_0|^2+|\alpha_1|^2=1
\]

\bigskip
the coefficients $\alpha_0,\alpha_1\in\mathbb{C}$ are called \textcolor{blue}{probability amplitudes}
\end{itemize}

\end{frame}

%%%%%%%%%%%%%%%%%%%%%%%%%%%%%%%%%%%%%%%%%%%%%%%%%%%%%%%%%%%%%%%

%% Quantum Register

%%%%%%%%%%%%%%%%%%%%%%%%%%%%%%%%%%%%%%%%%%%%%%%%%%%%%%%%%%%%%%%

\begin{frame}{Quantum register}

Quantum register: $2^n$-dimensional complex Hilbert space $(\C^2)^{\otimes n}$ with \textcolor{blue}{tensor product structure}

\[
(\C^2)^{\otimes n} := \underbrace{\mathbb{C}^2\otimes\mathbb{C}^2\otimes\cdots\otimes\mathbb{C}^2}_{n}
\]

\medskip
\begin{itemize}
\item Computational basis states (classical states):

\medskip
\[
|x\> = |x_1\> \otimes |x_2\> \otimes \cdots \otimes |x_n\>\,,\quad x\in\mathbb{B}^n
\]

\medskip
\item General state:

\medskip
\[
|\psi\>=\sum_{x\in\mathbb{B}^n} \alpha_x |x\>\,,\quad \sum_{x} |\alpha_x|^2 = 1
\]

\bigskip

\bigskip
Remark: Note that $|\psi\>$ is a column vector (ket) that is normalized with respect to the $\ell_2$-norm (Euclidean norm)
\end{itemize}
\end{frame}

%%%%%%%%%%%%%%%%%%%%%%%%%%%%%%%%%%%%%%%%%%%%%%%%%%%%%%%%%%%%%%%

%% Quantum Transformations

%%%%%%%%%%%%%%%%%%%%%%%%%%%%%%%%%%%%%%%%%%%%%%%%%%%%%%%%%%%%%%%

\begin{frame}{Quantum transformations}

Transformations on the quantum register $\cH:=(\C^2)^{\otimes n}$ are
described by \textcolor{blue}{unitary matrices} $U\in {\cal U}(\cH)$

\bigskip
Unitary matrices preserve the $\ell_2$-norm

\bigskip
Let $|\psi\>\in\cH$ be the state of the quantum register;  the state after the transformation $U$ is given by the matrix-vector product

\medskip
\begin{center}
{\large $U |\psi\>$}
\end{center}

\end{frame}

%%%%%%%%%%%%%%%%%%%%%%%%%%%%%%%%%%%%%%%%%%%%%%%%%%%%%%%%%%%%%%%

%% Quantum Circuits

%%%%%%%%%%%%%%%%%%%%%%%%%%%%%%%%%%%%%%%%%%%%%%%%%%%%%%%%%%%%%%%

\begin{frame}{Quantum circuit}

Each transformations $U$ has to be implemented by a \textcolor{red}{quantum circuit}, i.e., a sequence of
\textcolor{red}{elementary gates}

\bigskip
Quantum circuit model = Quantum mechanics + Notion of complexity

\end{frame}

%%%%%%%%%%%%%%%%%%%%%%%%%%%%%%%%%%%%%%%%%%%%%%%%%%%%%%%%%%%%%%%

%% Single Qubit Gate

%%%%%%%%%%%%%%%%%%%%%%%%%%%%%%%%%%%%%%%%%%%%%%%%%%%%%%%%%%%%%%%

\begin{frame}{Single qubit gate on two qubits}

single-qubit gate $U$ on \textcolor{red}{first} qubit\quad
\raisebox{0.5cm}{
 \Qcircuit @C=1em @R=1em {
  & \gate{{U}} & \qw \\
  & \qw      & \qw 
 }}

\pause

\bigskip
action on basis states of $\mathbb{C}^2\otimes\mathbb{C}^2$\quad
$
\begin{array}{ccc}
|0\> \otimes |0\> & \mapsto & ({\color{red} U}|0\>) \otimes |0\>\\
|0\> \otimes |1\> & \mapsto & ({\color{red} U}|0\>) \otimes |1\>\\
|1\> \otimes |0\> & \mapsto & ({\color{red} U}|1\>) \otimes |0\>\\
|1\> \otimes |1\> & \mapsto & ({\color{red} U}|1\>) \otimes |1\>
\end{array}
$

\pause

\bigskip
corresponding matrix 

${\color{red} U}\otimes {\color{blue} I}
=
\left(
\begin{array}{c|c}
{\color{red} u_{00}} \cdot {\color{blue} I} & {\color{red} u_{01}} \cdot {\color{blue} I} \\ \hline
{\color{red} u_{10}} \cdot {\color{blue} I} & {\color{red} u_{11}} \cdot {\color{blue} I}
\end{array}
\right)
=
\left(
\begin{array}{cccc}
{\color{red} u_{00}} & 0             & {\color{red} u_{01}} & 0      \\
0             & {\color{red} u_{00}} & 0             & {\color{red} u_{01}} \\
{\color{red} u_{10}} & 0             & {\color{red} u_{11}} & 0      \\
0             & {\color{red} u_{10}} & 0             & {\color{red} u_{11}}
\end{array}
\right)
$

\end{frame}

%%%%%%%%%%%%%%%%%%%%%%%%%%%%%%%%%%%%%%%%%%%%%%%%%%%%%%%%%%%%%%%

%% Single Qubit Gate

%%%%%%%%%%%%%%%%%%%%%%%%%%%%%%%%%%%%%%%%%%%%%%%%%%%%%%%%%%%%%%%

\begin{frame}{Single qubit gate on two qubits}

single-qubit gate $U$ on \textcolor{blue}{second} qubit\quad
\raisebox{0.5cm}{
 \Qcircuit @C=1em @R=1em {
  & \qw          & \qw \\
  & \gate{{U}}   & \qw 
 }}

\pause

\bigskip
action on basis states of $\mathbb{C}^2\otimes\mathbb{C}^2$\quad
$
\begin{array}{ccc}
|0\> \otimes |0\> & \mapsto & |0\> \otimes {\color{blue} U}|0\>\\
|0\> \otimes |1\> & \mapsto & |0\> \otimes {\color{blue} U}|1\>\\
|1\> \otimes |0\> & \mapsto & |1\> \otimes {\color{blue} U}|0\>\\
|1\> \otimes |1\> & \mapsto & |1\> \otimes {\color{blue} U}|1\>
\end{array}
$

\pause

\bigskip
corresponding matrix

$ {\color{red} I} \otimes {\color{blue} U}
=
\left(
\begin{array}{c|c}
{\color{red} 1} \cdot {\color{blue} U} & {\color{red} 0} \cdot {\color{blue} U} \\ \hline
{\color{red} 0} \cdot {\color{blue} U} & {\color{red} 1} \cdot {\color{blue} U}
\end{array}
\right)
=
\left(
\begin{array}{cccc}
{\color{blue} u_{00}} & {\color{blue} u_{01}} & 0      & 0 \\ 
{\color{blue} u_{10}} & {\color{blue} u_{11}} & 0      & 0 \\
0              & 0              & {\color{blue} u_{00}} & {\color{blue} u_{01}} \\
0              & 0              & {\color{blue} u_{10}} & {\color{blue} u_{11}}
\end{array}
\right)
$
\end{frame}


%%%%%%%%%%%%%%%%%%%%%%%%%%%%%%%%%%%%%%%%%%%%%%%%%%%%%%%%%%%%%%%

%% Controlled NOT gate

%%%%%%%%%%%%%%%%%%%%%%%%%%%%%%%%%%%%%%%%%%%%%%%%%%%%%%%%%%%%%%%

\begin{frame}{Controlled-NOT gate}

control: first qubit; \quad target: second qubit\quad
\raisebox{0.5cm}{
\Qcircuit @C=1em @R=1em {
& \ctrl{1} & \qw \\
& \targ     & \qw
}}

\bigskip
action on basis states of $\mathbb{C}^2\otimes\mathbb{C}^2$

\[
	|c\> \otimes |t\> \mapsto |c\> \otimes |c\oplus t\>
\]

\bigskip
corresponding matrix
\[
\left(
\begin{array}{cc|cc}
1      & 0      & 0      & 0 \\
0      & 1      & 0      & 0 \\ \hline
0      & 0      & 0      & {\color{blue}1} \\
0      & 0      & {\color{blue}1}      & 0
\end{array}
\right)
=
|0\>\<0| \otimes I_2 + |1\>\<1| \otimes {\color{blue} X}
\]

\medskip
where $I_2 = |0\>\<0|+|1\>\<1|$ and ${\color{blue}X=|0\>\<1| + |1\>\<0|}$
\end{frame}

%%%%%%%%%%%%%%%%%%%%%%%%%%%%%%%%%%%%%%%%%%%%%%%%%%%%%%%%%%%%%%%

%% Controlled U gate

%%%%%%%%%%%%%%%%%%%%%%%%%%%%%%%%%%%%%%%%%%%%%%%%%%%%%%%%%%%%%%%

\begin{frame}{Controlled $U$ gate}

control: qubit; \quad target: $m$-qubit register
\quad\raisebox{0.5cm}{
\Qcircuit @C=1em @R=1em {
& \qw & \ctrl{1} & \qw & \qw \\
& {/}\qw & \gate{U}     & \qw & \qw
}}

\bigskip
let $U$ be a unitary acting on the $m$-qubit register

\bigskip
action on basis states of $\mathbb{C}^2\otimes(\mathbb{C}^2)^{\otimes m}$

\[
	|c\> \otimes |t\> \mapsto |c\> \otimes U^c |t\> \quad \mbox{where } b\in\mathbb{B}, t\in\mathbb{B}^m
\]

\bigskip
corresponding matrix
\[
\left(
\begin{array}{c|c}
I & 0 \\ \hline
0 & {\color{blue}U} 
\end{array}
\right)
=
|0\>\<0| \otimes I + |1\>\<1| \otimes {\color{blue}U}
\]

\end{frame}

%%%%%%%%%%%%%%%%%%%%%%%%%%%%%%%%%%%%%%%%%%%%%%%%%%%%%%%%%%%%%%%

%% Toffoli gate

%%%%%%%%%%%%%%%%%%%%%%%%%%%%%%%%%%%%%%%%%%%%%%%%%%%%%%%%%%%%%%%

\begin{frame}{Toffoli gate}

control: first and second qubits; \quad target: third qubit\quad
\raisebox{0.5cm}{
\Qcircuit @C=1em @R=1em {
& \ctrl{2} & \qw \\
& \ctrl{1} & \qw \\
& \targ     & \qw
}}

\bigskip
action on basis states of $\C^2\otimes\C^2\otimes\C^2$

\[
|c_1\> \otimes |c_2\> \otimes |t\> \mapsto |c_1\> \otimes |c_2\> \otimes |(c_1\land c_2)\oplus t\>
\]

\bigskip
corresponding matrix
\[
\left(
\begin{array}{cc|cc|cc|cc}
1 & 0 & 0 & 0 & 0 & 0 & 0 & 0 \\ 
0 & 1 & 0 & 0 & 0 & 0 & 0 & 0 \\ \hline
0 & 0 & 1 & 0 & 0 & 0 & 0 & 0 \\ 
0 & 0 & 0 & 1 & 0 & 0 & 0 & 0 \\ \hline
0 & 0 & 0 & 0 & 1 & 0 & 0 & 0 \\ 
0 & 0 & 0 & 0 & 0 & 1 & 0 & 0 \\ \hline 
0 & 0 & 0 & 0 & 0 & 0 & 0 & {\color{blue}1} \\ 
0 & 0 & 0 & 0 & 0 & 0 & {\color{blue}1} & 0  
\end{array}
\right)
=
(I_4 - |11\>\<11|) \otimes I_2 + |11\>\<11| \otimes {\color{blue}X}
\]

\end{frame}

%%%%%%%%%%%%%%%%%%%%%%%%%%%%%%%%%%%%%%%%%%%%%%%%%%%%%%%%%%%%%%%

%% Reversible Computation

%%%%%%%%%%%%%%%%%%%%%%%%%%%%%%%%%%%%%%%%%%%%%%%%%%%%%%%%%%%%%%%

\begin{frame}{Simulating irreversible gates with Toffoli gate}

The classical AND gate is irreversible because if the output is $0$ then we cannot determine which of the three possible pairs was the actual input

\bigskip

\hspace{1cm}
\begin{tabular}{cc|c}
$x_1$ & $x_2$ & $x_1 \land x_2$ \\ \hline
0 & 0 & 0 \\
0 & 1 & 0 \\
1 & 0 & 0 \\
1 & 1 & 1 
\end{tabular}

\pause

\bigskip
But it is easy to simulate the AND gate with one Toffoli gate

\hspace{1cm}
\Qcircuit @C=1em @R=1em {
\lstick{|x_1\>} & \qw & \ctrl{2} & \qw & \rstick{|x_1\>} \\
\lstick{|x_2\>} & \qw & \ctrl{1} & \qw & \rstick{|x_2\>} \\
\lstick{|0\>}   & \qw & \targ    & \qw & \rstick{|x_1 \land x_2\>}
}

\end{frame}

%%%%%%%%%%%%%%%%%%%%%%%%%%%%%%%%%%%%%%%%%%%%%%%%%%%%%%%%%%%%%%%

%% Garbage collection

%%%%%%%%%%%%%%%%%%%%%%%%%%%%%%%%%%%%%%%%%%%%%%%%%%%%%%%%%%%%%%%

\begin{frame}{Problem of garbage}

To simulate irreversible circuits with Toffoli gates, we keep the input and intermediary results to make everything reversible  

\bigskip
Consider the function $y=x_1\land x_2\land x_3$

\bigskip
\hspace{1cm}
\Qcircuit @C=1em @R=1em {
\lstick{|x_1\>} & \qw & \ctrl{4} & \qw      & \qw & \rstick{|x_1\>} \\
\lstick{|x_2\>} & \qw & \ctrl{3} & \qw      & \qw & \rstick{|x_2\>} \\
\lstick{|x_3\>} & \qw & \qw      & \ctrl{1} & \qw & \rstick{|x_3\>} \\
\lstick{|0\>}   & \qw & \qw      & \targ    & \qw & \rstick{|x_1 \land x_2 \land x_3\>} \\
\lstick{|0\>}   & \qw & \targ    & \ctrl{-1}& \qw & \rstick{|x_1 \land x_2\> \mbox{ garbage}}
}

\pause

\bigskip

\bigskip
It is important to not leave any garbage;  otherwise, we could not make use of quantum parallelism and constructive interference effects

\end{frame}

%%%%%%%%%%%%%%%%%%%%%%%%%%%%%%%%%%%%%%%%%%%%%%%%%%%%%%%%%%%%%%%

%% Reversible Computation

%%%%%%%%%%%%%%%%%%%%%%%%%%%%%%%%%%%%%%%%%%%%%%%%%%%%%%%%%%%%%%%

\begin{frame}{Reversible garbage removal}

It is always possible to reversibly remove (uncompute) the garbage

\bigskip
In the case $y=x_1 \land x_2 \land x_3$, this can be done with the circuit


\bigskip
\hspace{1cm}
\Qcircuit @C=1em @R=1em {
\lstick{|x_1\>} & \qw & \ctrl{4} & \qw      & \ctrl{4} & \qw & \rstick{|x_1\>} \\
\lstick{|x_2\>} & \qw & \ctrl{3} & \qw      & \ctrl{3} & \qw & \rstick{|x_2\>} \\
\lstick{|x_3\>} & \qw & \qw      & \ctrl{1} & \qw      & \qw & \rstick{|x_3\>} \\
\lstick{|0\>}   & \qw & \qw      & \targ    & \qw      & \qw & \rstick{|x_1 \land x_2 \land x_3\>} \\
\lstick{|0\>}   & \qw & \targ    & \ctrl{-1}& \targ    & \qw & \rstick{|0\> \mbox{ garbage uncomputed}}
}


\end{frame}

%%%%%%%%%%%%%%%%%%%%%%%%%%%%%%%%%%%%%%%%%%%%%%%%%%%%%%%%%%%%%%%

%% Reversible Computation

%%%%%%%%%%%%%%%%%%%%%%%%%%%%%%%%%%%%%%%%%%%%%%%%%%%%%%%%%%%%%%%

\begin{frame}{Simulating irreversible circuits with Toffoli gates}

Let $f:\{0,1\}^n\rightarrow \{0,1\}$ be any boolean function \pause

\bigskip
Assume this function can be computed classically using only $t$ classical elementary gates such as AND, OR, NAND \pause

\bigskip
We can implement a unitary $U_f$ on $(\C^2)^{\otimes n} \otimes \C^2 \otimes (\C^2)^{\otimes w}$ such that

\[
U_f \big( |x\>_{\mathrm{in}} \otimes |y\>_{\mathrm{out}} \otimes |0\>^{\otimes w}_{\mathrm{work}} \big) = |x\> \otimes |y\oplus f(x) \> \otimes |0\>^{\otimes w}
\]

\pause

\bigskip
$U_f$ is built from polynomially many in $t$ Toffoli gates and the size $w$ of the workspace register is polynomial in $t$

\pause

\bigskip
During the computation the qubits of the workspace register are changed, but at the end they reversibly reset to $|0\>^{\otimes w}$

\end{frame}


%%%%%%%%%%%%%%%%%%%%%%%%%%%%%%%%%%%%%%%%%%%%%%%%%%%%%%%%%%%%%%%

%% Universal Gate Set -- Exact Implementation

%%%%%%%%%%%%%%%%%%%%%%%%%%%%%%%%%%%%%%%%%%%%%%%%%%%%%%%%%%%%%%%

\begin{frame}{Universal gate set -- exact implementation}

Each unitary $U\in {\cal U}(\cH)$ can be
implemented exactly by quantum circuits using only:

\medskip
\begin{itemize}
\item \textcolor{blue}{CNOT} gates (acting on adjacent qubits)

\bigskip
\item arbitrary \textcolor{blue}{single qubit} gates
\end{itemize}

\end{frame}

%%%%%%%%%%%%%%%%%%%%%%%%%%%%%%%%%%%%%%%%%%%%%%%%%%%%%%%%%%%%%%%

%% Gate Complexity of Unitaries -- Exact Implementation

%%%%%%%%%%%%%%%%%%%%%%%%%%%%%%%%%%%%%%%%%%%%%%%%%%%%%%%%%%%%%%%

\begin{frame}{Gate complexity of unitaries -- exact implementation}

The gate complexity $\kappa(U)$ of a unitary $U\in {\cal U}(\cH)$ is minimal number of elementary gates needed to implement $U$

\bigskip

\bigskip
For example, quantum Fourier Transform has complexity $O({\color{red} n^2})$ 

\medskip
\quad $\Longrightarrow$\quad Shor's factorization algorithm

\end{frame}

%%%%%%%%%%%%%%%%%%%%%%%%%%%%%%%%%%%%%%%%%%%%%%%%%%%%%%%%%%%%%%%

%% Universal Gate Set -- Approximate Implementation

%%%%%%%%%%%%%%%%%%%%%%%%%%%%%%%%%%%%%%%%%%%%%%%%%%%%%%%%%%%%%%%

\begin{frame}{Universal gate set -- approximate implementation}

For each $\epsilon\in (0,1)$ and each unitary $U\in {\cal U}(\cH)$, there is a unitary $V$ such that

\[
\| U - V \| \le \epsilon \quad\mbox{where}\quad \|U-V\| = \sup_{|\psi\>} \|(U-V)|\psi\>\|
\]

and $V$ is implemented by quantum circuits using only:

\medskip
\begin{itemize}
\item \textcolor{blue}{CNOT} gates (acting on adjacent qubits)

\bigskip
\item the \textcolor{blue}{single qubit} gates

\[
H = \nicefrac{1}{\sqrt{2}}
\left( 
\begin{array}{cc}
1 & 1 \\
1 & -1
\end{array}
\right)\,\quad
R(\theta) = 
\left( 
\begin{array}{cc}
1 & 0 \\
0 & e^{i\theta}
\end{array}
\right)\,,
\quad\mbox{with } \theta=\frac{\pi}{4}
\]
\end{itemize}

\pause

\bigskip
There are other universal gate sets

\end{frame}

%%%%%%%%%%%%%%%%%%%%%%%%%%%%%%%%%%%%%%%%%%%%%%%%%%%%%%%%%%%%%%%

%% Gate Complexity of Unitaries -- Exact Implementation

%%%%%%%%%%%%%%%%%%%%%%%%%%%%%%%%%%%%%%%%%%%%%%%%%%%%%%%%%%%%%%%

\begin{frame}{Gate complexity of unitaries -- approximate implementation}

The gate complexity $\kappa_\epsilon(U)$ of a unitary $U$ is the minimal number of gates (from a universal gate set) need to implement a unitary $V$ with $\|U-V\|\le\epsilon$

\pause

\bigskip

\bigskip
The Solovay-Kitaev theorem implies that 
\[
\kappa_\epsilon(U) = O\Big(\kappa(U) \cdot \log^c\big(\kappa(U)/\epsilon\big)\Big)
\]
for some small constant $c$

\pause

\bigskip

\bigskip
Counting arguments show that most $n$-qubit unitaries have gate complexity \textcolor{red}{exponential} in $n$.

\end{frame}

%%%%%%%%%%%%%%%%%%%%%%%%%%%%%%%%%%%%%%%%%%%%%%%%%%%%%%%%%%%%%%%

%% Quantum Measurement

%%%%%%%%%%%%%%%%%%%%%%%%%%%%%%%%%%%%%%%%%%%%%%%%%%%%%%%%%%%%%%%

\begin{frame}{Quantum measurement}

A general measurement is described by a collection $P_0$,\ldots,$P_{m-1}$ of orthogonal projectors such that 

\[
\sum_{i=0}^{m-1} P_i = I_{\cH}\quad\mbox{where $\cH$ denotes the identity on $\cH$}
\]

\pause

\bigskip

Let $|\psi\>$ be the state of the quantum register.  The probability of obtaining the outcome $i$ is given by
\[
\Pr(i) = \| P_i |\psi\> \|^2
\]

\bigskip

\pause
The post-measurement state (collapse of the wavefunction) is 
\[
\frac{P_i|\psi\>}{\|P_i |\psi\>\|}
\]

\end{frame}

%%%%%%%%%%%%%%%%%%%%%%%%%%%%%%%%%%%%%%%%%%%%%%%%%%%%%%%%%%%%%%%

%% Quantum Measurements

%%%%%%%%%%%%%%%%%%%%%%%%%%%%%%%%%%%%%%%%%%%%%%%%%%%%%%%%%%%%%%%

\begin{frame}{Elementary quantum measurements}

A measurement has to be realized by first applying a suitable quantum circuit followed by an elementary measurement

\pause

\bigskip

An elementary measurement on the $n$-qubit quantum register $\cH$ consists of measuring the first (w.l.o.g.) $m$ qubits $(m\le n$) with respect to the computational basis

\pause

\bigskip
The $2^m$ orthogonal projectors $P_b$ are labeled by $m$-bit strings $b\in\mathbb{B}^m$ and are defined by
\[
P_b = |b_1\>\<b_1| \otimes |b_2\>\<b_2| \otimes \cdots \otimes |b_m\>\<b_m| \otimes I_{2^{n-m}}
\]

\pause

\bigskip
The probability of obtaining outcome $b$ is given by
\[
\Pr(b) = \| P_b |\psi\> \|^2 = \sum_{x_{m+1},\ldots,x_n\in\mathbb{B}} |\alpha_{b_1,\ldots,b_m,x_{m+1},\ldots,x_n}|^2
\]

\end{frame}

%%%%%%%%%%%%%%%%%%%%%%%%%%%%%%%%%%%%%%%%%%%%%%%%%%%%%%%%%%%%%%%

%% Quantum algorithms

%%%%%%%%%%%%%%%%%%%%%%%%%%%%%%%%%%%%%%%%%%%%%%%%%%%%%%%%%%%%%%%

\begin{frame}{Structure of quantum algorithms}

A quantum algorithm consists of 

\bigskip
\begin{itemize}
\item preparing the initial state $|x\>$ with $x\in\mathbb{B}^n$,

\bigskip
\item applying a quantum circuit of \textcolor{blue}{polynomially} many in $n$ gates from some universal gate set, and

\bigskip
\item performing an elementary measurement
\end{itemize}

\bigskip
These steps are repeated polynomially many times to collect enough samples and followed by classical post-processing $\Longrightarrow$ solution of the problem


\end{frame}

%%%%%%%%%%%%%%%%%%%%%%%%%%%%%%%%%%%%%%%%%%%%%%%%%%%%%%%%%%%%%%%

\begin{frame}{Hadamard test}

\hspace{4cm}
\Qcircuit @C=1em @R=0.7em {
\lstick{|0\>}    & \gate{H}    & \ctrl{1} & \gate{H} & \meter \\
\lstick{|\psi\>} & {/}\qw      & \gate{U} & \qw      & \qw
}

\pause

\bigskip

The probabilities of obtaining the outcomes $0$ and $1$ are:
\[
\Pr(0) = \nicefrac{1}{2}\left( 1 + \mathrm{Re}\<\psi|U|\psi\> \right) \quad\quad
\Pr(1) = \nicefrac{1}{2}\left( 1 - \mathrm{Re}\<\psi|U|\psi\> \right) 
\]

\end{frame}

%%%%%%%%%%%%%%%%%%%%%%%%%%%%%%%%%%%%%%%%%%%%%%%%%%%%%%%%%%%%%%%

%% Hadamard test

%%%%%%%%%%%%%%%%%%%%%%%%%%%%%%%%%%%%%%%%%%%%%%%%%%%%%%%%%%%%%%%

\begin{frame}{Hadamard test}

\hspace{4cm}
\Qcircuit @C=1em @R=0.7em {
\lstick{|0\>}    & \gate{H}    & \ctrl{1} & \gate{H} & \meter \\
\lstick{|\psi\>} & {/}\qw      & \gate{U} & \qw      & \qw
}

\pause

\bigskip

\begin{eqnarray*}
& & 
|0\> \otimes |\psi\> \\ \pause
& \mapsto &
\nicefrac{1}{\sqrt{2}}\, (|0\> + |1\>) \otimes |\psi\> \\ \pause
& = &
\nicefrac{1}{\sqrt{2}} |0\> \otimes |\psi\> + \nicefrac{1}{\sqrt{2}} |1\> \otimes |\psi\> \\ \pause
& \mapsto &
\nicefrac{1}{\sqrt{2}} |0\> \otimes |\psi\> + \nicefrac{1}{\sqrt{2}} |1\> \otimes U |\psi\> \\ \pause
& \mapsto &
\nicefrac{1}{2} (|0\> + |1\>) \otimes |\psi\> + \nicefrac{1}{2} (|0\> - |1\>) \otimes U |\psi\> \\ \pause
& = &
\nicefrac{1}{2} |0\> \otimes (|\psi\> + U|\psi\>) + \nicefrac{1}{2} |1\> \otimes (|\psi\> - U |\psi\>) \\
& =: & |\Phi\>
\end{eqnarray*}

\end{frame}

%%%%%%%%%%%%%%%%%%%%%%%%%%%%%%%%%%%%%%%%%%%%%%%%%%%%%%%%%%%%%%%

%% Hadamard test

%%%%%%%%%%%%%%%%%%%%%%%%%%%%%%%%%%%%%%%%%%%%%%%%%%%%%%%%%%%%%%%

\begin{frame}{Hadamard test}

\begin{eqnarray*}
& & 
\Pr(0) \\
& = & 
\| P |\Phi\> \|^2 \quad \mbox{with} \quad P = |0\>\<0| \otimes I \\ \pause
& = &
\| \big( {\color{red} |0\>\<0|} \otimes I \big) \, \Big( \nicefrac{1}{2} {\color{red} |0\>} \otimes (|\psi\> + U|\psi\>) + \nicefrac{1}{2} {\color{blue} |1\>} \otimes (|\psi\> - U |\psi\>) \Big) \|^2 \\ \pause
& = & 
\| \nicefrac{1}{2} |0\> \otimes (|\psi\> + U|\psi\>) \|^2 \\ \pause
& = &
\nicefrac{1}{4} \, \| |0\> \|^2 \cdot \| |\psi\> + U|\psi\> \|^2 \\ \pause
& = &
\nicefrac{1}{4} \, \big( \<\psi| + \<\psi|U^\dagger\big) \, \big( |\psi\> + U|\psi\> \big) \\ \pause
& = &
\nicefrac{1}{4} \, \big( \<\psi|\psi\> + \<\psi|U|\psi\> + \<\psi|U^\dagger|\psi\> + \<\psi|U^\dagger U|\psi\>\big) \\ \pause
& = & 
\nicefrac{1}{4} \, \big( 2 + \<\psi|U|\psi\> + \overline{\<\psi|U|\psi\>}\big) \\ \pause
& = &
\nicefrac{1}{2} \, \big( 1 + \mathrm{Re}\<\psi|U|\psi\> \big) \\
\end{eqnarray*}

\end{frame}

%%%%%%%%%%%%%%%%%%%%%%%%%%%%%%%%%%%%%%%%%%%%%%%%%%%%%%%%%%%%%%%

%% Hadamard test

%%%%%%%%%%%%%%%%%%%%%%%%%%%%%%%%%%%%%%%%%%%%%%%%%%%%%%%%%%%%%%%

\begin{frame}{Hadamard test -- Figure it out yourself}

How can you estimate the imaginary part of $\<\psi|U|\psi\>$?

\bigskip

Hint: Add a simple gate on the control register before the measurement.

\end{frame}

%%%%%%%%%%%%%%%%%%%%%%%%%%%%%%%%%%%%%%%%%%%%%%%%%%%%%%%%%%%%%%%

%% SWAP Test

%%%%%%%%%%%%%%%%%%%%%%%%%%%%%%%%%%%%%%%%%%%%%%%%%%%%%%%%%%%%%%%

\begin{frame}{SWAP test -- Figure it out yourself}

Let $S$ denote the swap gate acting on two qubits
\[
S = |00\>\<00| + |01\>\<10| + |10\>\<01| + |11\>\<11|
\]

\bigskip
Determine the matrix representation of $S$ with respect to the computational basis

\bigskip
Consider the Hadamard test where the controlled operation is
\[
|0\>\<0|\otimes I_4 + |1\>\<1| \otimes S
\]
and the state of the target register $|\psi\>=|\psi_1\>\otimes |\psi_2\>$

\bigskip
Determine the probability of obtaining $0$ and $1$ for the cases:
\begin{itemize}
\item arbitrary $|\psi_1\>$ and $|\psi_2\>$, 
\item $\<\psi_1|\psi_2\>=0$ (orthogonal), and 
\item $\<\psi_1|\psi_2\>=1$ (the same).
\end{itemize}
\end{frame}

